\section*{第 \ref{ch:b1:chemical-biogenic-sediments} 章——碳酸盐、硅质与蒸发岩沉积}

\begin{solution}{exo:b1:chemical-biogenic-sediments:1}
氯 $0.5503\times35.165 = \qty{19.35}{g}$；钠 $0.3066\times35.165 =
\qty{10.78}{g}$；硫酸盐 \qty{2.711}{g}；镁 \qty{1.284}{g}。钠与
氯合计 $0.5503 + 0.3066 = 0.857$，占
盐分的 \qty{85.7}{\%}。
\end{solution}

\begin{solution}{exo:b1:chemical-biogenic-sediments:2}
剩余分数为 $1/\mathrm{CF}$：碳酸钙 $1/1.8 = 0.556$ 剩余，
\qty{44.4}{\%} 被蒸发；石膏 $0.263$ 剩余，\qty{73.7}{\%} 被蒸发；
岩盐 $0.094$ 剩余，\qty{90.6}{\%} 被蒸发；硫酸镁 $0.014$ 剩余，
\qty{98.6}{\%} 被蒸发。
\end{solution}

\begin{solution}{exo:b1:chemical-biogenic-sediments:3}
(a)~$\rho_s(1-\phi) = G/a = 8.4/0.0042 = \qty{2000}{kg.m^{-3}}$。(b)~取文石时，
$1 - 2000/2930 = 0.32$，孔隙度为 \qty{32}{\%}；取方解石时，
$1 - 2000/2710 = 0.26$。(c)~$a = 3.7/2000 = \qty{1.85e-3}{m}$，即每
年 \qty{1.85}{mm}。
\end{solution}

\begin{solution}{exo:b1:chemical-biogenic-sediments:4}
在 \qty{4000}{m} 以浅，含量下降缓慢（由 88 降至 80）；
在 \qty{4000}{m} 以深则急剧下降，\omterm{def:b1:chemical-biogenic-sediments:ccd}{溶跃面}约位于 \qty{4000}{m} 处。
含量在 \qty{4700}{m}（12）与 \qty{4800}{m}
（6）之间降至 \qty{10}{\%}，$4700 + 100\times(12-10)/(12-6) = \qty{4733}{m}$，即
\omterm{def:b1:chemical-biogenic-sediments:ccd}{碳酸盐补偿深度}约为 \qty{4730}{m}，比\omterm{def:b1:chemical-biogenic-sediments:ccd}{溶跃面}深约 \qty{700}{m}。
\end{solution}

\begin{solution}{exo:b1:chemical-biogenic-sediments:5}
(a)~\ce{Ca^2+ + 2 HCO3- -> CaCO3 + CO2 + H2O}：每摩尔碳酸钙生成 1 摩尔 \ce{CO2}，
故每千克 \ce{CaCO3} 对应 $44.01/100.09 = \qty{0.440}{kg}$ 的 \ce{CO2}，
(b)~每平方米每年 \ce{CaCO3} 对应 $3.7\times0.440 = \qty{1.63}{kg}$ 的 \ce{CO2}。
(c)~用掉的两个碳酸氢根离子中，只有一个碳进入
\omterm{def:b1:rock-forming-minerals:mineral}{矿物}，另一个则以二氧化碳形式逸出。钙化作用消耗海水中的碱度，
并降低水体容纳溶解碳的能力；其中有多少
\ce{CO2} 进入大气，取决于海水的缓冲作用，
这属于化学教材中碳酸盐平衡的内容。
\end{solution}

\begin{solution}{exo:b1:chemical-biogenic-sediments:6}
(a)~$t_c = (2000/350)^2 = \qty{32.7}{Ma}$；按每百万年 \qty{25}{km} 计，
$25\times32.7 = \qty{816}{km}$，距轴线。(b)~$t_c = (2600/350)^2 =
\qty{55.2}{Ma}$，即 \qty{1380}{km}。(c)~\qty{60}{Ma} 龄的洋底已接收
红黏土的时间：第一种情形下为 $60 - 32.7 = \qty{27.3}{Ma}$，而
第二种情形下仅为 $60 - 55.2 = \qty{4.8}{Ma}$。
\end{solution}

\begin{solution}{exo:b1:chemical-biogenic-sediments:7}
(a)~是的：$a = 4.2 > r = 3$，礁体相对海平面每年增长 \qty{1.2}{mm}，
直至礁顶到达水面。(b)~$D$ 以每年 $10 - 4.2 = \qty{5.8}{mm}$ 增长，
由 2 增至 \qty{50}{m} 需 $48/0.0058 = \num{8300}$ 年。(c)~因为
其生长速率随水深增加而下降：到达藻类的光线减少，
这些藻类生活在珊瑚体内并为其提供营养，因此 $a$ 减小，礁体
相对海平面落后得越来越快。
\end{solution}

\begin{solution}{exo:b1:chemical-biogenic-sediments:8}
(a)~反应为 \ce{CaSO4.2H2O -> CaSO4 + 2 H2O}。(b)~摩尔体积分别为 $172.17/2.30 =
\qty{74.9}{cm^3}$ 与 $136.14/2.97 = \qty{45.8}{cm^3}$，因此固体体积减少
$1 - 45.8/74.9 = \qty{38.8}{\%}$。(c)~$2\times18.02/172.17 = 0.209$：
\qty{20.9}{\%} 的质量以水的形式离开。这些水在深部释放后，
会提高周围岩石孔隙中的压强；若被盖层封住，它会使
岩石变弱并助长断层滑动，还可溶解并携带其他
盐类。
\end{solution}

\begin{solution}{exo:b1:chemical-biogenic-sediments:9}
(a)~$9.2/255 = \qty{3.6}{\%}$ 被保存，\qty{96.4}{\%} 被溶解。(b)~在
生产力很高时，即便只保留其中百分之几，仍是很大的通量；而在
很少有黏土或碳酸盐稀释它的地区：南大洋、北太平洋、
赤道太平洋的上升流区，以及\omterm{def:b1:chemical-biogenic-sediments:ccd}{碳酸盐补偿深度}以下。
(c)~薄而精细的硅藻壳最先溶解，坚固者得以保存：
沉积物保存的是最坚固物种的一部分，而非
水中曾经生存的群落。
\end{solution}

\begin{solution}{exo:b1:chemical-biogenic-sediments:10}
(a)~$EA = 1.5\times\qty{1e10}{m^2} = \qty{1.5e10}{m^3}$ 每年。(b)~当
$Q_{\mathrm{in}} = \qty{1.8e10}{m^3}$，$Q_{\mathrm{out}} = \qty{0.3e10}{m^3}$ 且
$k = 6$：超过石膏阈值（3.8），但未达岩盐阈值（10.6），因此
碳酸钙与石膏沉淀。当 $Q_{\mathrm{in}} =
\qty{1.575e10}{m^3}$，$Q_{\mathrm{out}} = \qty{0.075e10}{m^3}$ 且 $k = 21$ 时，对于
镁这类留在溶液中的离子：卤水超过岩盐阈值，
因此岩盐也沉淀，但未达镁盐的阈值（70）。
(c)~岩盐沉淀要求 $Q_{\mathrm{in}}/(Q_{\mathrm{in}} - EA) \ge 10.6$，即
$Q_{\mathrm{in}} \le 10.6\,EA/9.6 = 1.104\,EA = \qty{1.66e10}{m^3}$ 每年：
入流量至多超过蒸发量约十分之一。
\end{solution}

\begin{solution}{exo:b1:chemical-biogenic-sediments:11}
(a)~海水每千克含溴化物 \qty{67.3}{mg}；浓缩至 10.6~倍时，
卤水含 $67.3\times10.6 = \qty{713}{mg.kg^{-1}}$，首批岩盐
含 $0.032\times713 = \qty{22.8}{mg.kg^{-1}}$。(b)~浓缩至 30~倍时，$67.3\times30 =
\qty{2019}{mg.kg^{-1}}$（卤水中），岩盐中为 \qty{64.6}{mg.kg^{-1}}。
(c)~岩盐只占卤水浓度的 3.2~个百分之一，因此
溴化物几乎全部留在溶液中，并随水分蒸发而持续浓缩。
(d)~上升记录了卤水的浓缩；骤降则记录了卤水被稀释，
由新的海水或淡水流入盆地所致，
发生于岩盐继续生长之前。
\end{solution}

\begin{solution}{exo:b1:chemical-biogenic-sediments:12}
(a)~$t_c = \qty{32.65}{Ma}$，故软泥厚 $0.02\times\num{32650} = \qty{653}{m}$
（\num{32650} 个千年间隔，每间隔 \qty{2}{cm}）；红黏土厚
$0.001\times\num{37350} = \qty{37}{m}$。(b)~$t_c = (2500/350)^2 =
\qty{51.0}{Ma}$：软泥厚 \qty{1020}{m}，红黏土厚 \qty{19}{m}。较深的\omterm{def:b1:chemical-biogenic-sediments:ccd}{碳酸盐补偿深度}
使软泥厚得多，而深大西洋洋底的红黏土
比同龄太平洋洋底的更薄。(c)~\omterm{def:b1:chemical-biogenic-sediments:ccd}{碳酸盐补偿深度}可能曾在
一段时间内更深，使该地点重新升至其上方；或
板块将该地点带入高生产力带之下（如赤道），那里有更多碳酸盐
沉降下来；或该层由滑坡或\omterm{def:b1:clastic-sediments:turbidite}{浊流}从
海底较高部位搬运而来。
\end{solution}

\begin{solution}{pb:b1:chemical-biogenic-sediments:1}
\textbf{1.} $1030\times1000 = \qty{1.03e6}{kg}$ 的海水，其中
含盐 $1.03\times10^6\times0.035165 = \qty{36200}{kg}$。

\textbf{2.} 石膏 $1.03\times10^6\times1.581\times10^{-3} = \qty{1628}{kg}$；
石盐 $1.03\times10^6\times0.02741 = \qty{28230}{kg}$。

\textbf{3.} 碳酸钙 $113.3/2710 = \qty{0.042}{m}$；石膏 $1628/2300 =
\qty{0.71}{m}$；石盐 $28230/2170 = \qty{13.0}{m}$。

\textbf{4.} 硬石膏的质量为 $1628\times136.14/172.17 =
\qty{1288}{kg}$，厚度为 $1288/2970 = \qty{0.43}{m}$。

\textbf{5.} 干盐的含量为 $0.110 + 1.250 + 27.41 + 6.34 = \qty{35.11}{g}$ 每
千克，其中石盐占 $27.41/35.11 = \qty{78}{\%}$：钠和氯
构成溶解盐的大部分，石盐几乎占去了它们的全部。

\textbf{6.} 自下而上依次为碳酸钙、石膏、石盐、苦卤
盐类。每种\omterm{def:b1:rock-forming-minerals:mineral}{矿物}在卤水达到其
饱和时开始沉淀，溶解度最小者先沉淀，并沉积于先前形成的层之上。

\textbf{7.} 卤水深度为 $1000/\mathrm{CF}$：碳酸钙 \qty{556}{m}，
石膏 \qty{263}{m}，石盐 \qty{94}{m}。

\textbf{8.} $1 - 1/10.6 = \qty{90.6}{\%}$。

\textbf{9.} 硫酸镁：$1000/70 = \qty{14}{m}$，钾盐：
$1000/90 = \qty{11}{m}$。

\textbf{10.} 碳酸钙形成时卤水密度约为 \qty{1.10}{g.cm^{-3}}，石膏时为 1.13，
石盐时为 1.22，苦卤时超过 1.30，而
海水为 1.03。流入的海水浮于致密卤水之上，卤水
留在底部；在有海槛的盆地中，卤水可以
从流入水下方流回大洋。

\textbf{11.} 其厚度仅 \qty{4}{cm}，而石盐厚 \qty{13}{m}；它
形成于边缘，那里流入的海水先行浓缩，浓缩程度高于中部，
因此它与石膏底部混杂，或被石膏所取代。

\textbf{12.} 石盐只夺取了卤水中溴化物的一小部分
（练习~\ref{exo:b1:chemical-biogenic-sediments:11}），因此卤水随石盐生长而变得
溴化物更加富集，每新生一层石盐便又夺取
更多。

\textbf{13.} $1500/13.01 = 115.3$：需 116 次充填。

\textbf{14.} 硬石膏共 $116\times0.434 = \qty{50}{m}$，不是单一层，而是
116 层，每层厚度略小于半米，位于各自充填期石盐的底部：
构成硬石膏与石盐交替的旋回层序。

\textbf{15.} 每次充填，$1.03\times10^6\times6.34\times10^{-3}/2000 =
\qty{3.3}{m}$，故 116 次充填共 \qty{379}{m}。它们是最易溶的
盐类：下一次海水流入会将其重新溶解；若盆地
卤水在演化终了前流出，则永远达不到这些盐的析出条件。

\textbf{16.} 每次再充填时卤水都被稀释：溴化物含量在
新石盐底部急剧下降，向上又回升，形成
锯齿状剖面，每次充填对应一个齿。

\textbf{17.} $2000/13.01 = 154$ 次充填。

\textbf{18.} $5000\times365.25\times86400 = \qty{1.578e11}{m^3}$，每年
均匀分布于 $\qty{1.0e11}{m^2}$ 上：$E = \qty{1.58}{m}$。

\textbf{19.} 每年的盐量为 $1.578\times10^{11}\times1030\times0.035165 = \qty{5.72e12}{kg}$
。

\textbf{20.} $1.578\times10^{11}\times1030\times0.02741/2170 =
\qty{2.05e9}{m^3}$ 的石盐，即 \qty{0.0205}{m}，约 \qty{2}{cm}
每年，堆积于盆底之上。

\textbf{21.} $1500/0.0205 = \num{73000}$ 年。

\textbf{22.} 水位以每年 \qty{1.58}{m} 的速率下降，\qty{1000}{m} 深的水
在 $1000/1.578 = 634$ 年内即可耗尽。

\textbf{23.} 当水量减少 \qty{90.6}{\%} 时石盐开始出现，
历时 $905.7/1.578 = 574$ 年。

\textbf{24.} \qty{1.5}{km} 厚的石盐，其盐分取自盆地的 116 次充填
（115.3，向上取整为 116），这些充填无论是随盆地
反复干涸与再充填而依次发生，还是如第~21~题所示，连续越过海槛，
以每
$1000/1.578 = 634$ 年补入一盆海水的速率进行。
\end{solution}
